% GENERATED FILE. Edit canonical lessons, then run npm run generate:books.
% canonical-source-hash: fnv1a64:4f37da60f2034ba4
% canonical-lessons: ZH-C532-youwang, ZH-W532-wang-hope, ZH-C532-babude, ZH-C532-qiwang, ZH-C532-ruyi, ZH-C532-zhiwang

\chapter{Hoping for More}
\label{ch:zh-hoping-for-more}
\hlchaptermodality{532}

\begin{chapteropening}
\textbf{By the end of this chapter:} \emph{I can say I'm dying for something, that something should happen, what I expect and whom I count on.}

Complete the last lesson of chapter 532: I can say I'm dying for something, that something should happen, what I expect and whom I count on.
\end{chapteropening}

\section[\texorpdfstring{\zh{有望} (yǒuwàng)}{有望 (yǒuwàng)}]{\zh{有望} (yǒuwàng) — with good hopes of it; likely to happen}

\label{lesson:ZH-C532-youwang}



\begin{quote}
Before the new one: say the Mandarin for a misfortune, then the Mandarin for to lose one's grip.
\end{quote}

\subsection*{You'll want to know: \zh{有望}}
\textbf{\zh{有望}} — \emph{yǒuwàng} — \textquotedblleft{}with good hopes of it; likely to happen\textquotedblright{}.

\begin{quote}
\textbf{\zh{如果天气好飞机就有望正点起飞}} \emph{rúguǒ tiānqì hǎo fēijī jiù yǒuwàng zhèngdiǎn qǐfēi} — If the weather is good, the plane should take off on time.
\end{quote}

The new character, in pieces: \zh{望} — eleven strokes in three pieces: on the top left a dot, a bar and a short vertical that turns right, three strokes; on the top right \zh{月}, four strokes; and \zh{王} at the foot, four strokes. The next lesson puts it in your hand.

\begin{cousinweb}
\textbf{\zh{有望}} (\emph{yǒuwàng}) is built from \zh{有} 'to have' and \emph{wàng}, both inherited from Middle Chinese, \emph{hjuwX} and \emph{mjangH} (Guangyun, in Baxter's transcription), which glosses the second 'to look out for' and, after an older dictionary, as watching for someone far away to come back.
\end{cousinweb}

\begin{grammarlens}[title={What you've built}]
One more word for saying what would happen.
\end{grammarlens}

\subsection*{Your turn}
\begin{itemize}
\raggedright
  \item \emph{Say it:} \emph{yǒuwàng}
  \item \emph{Say it:} \emph{yǒuwàng}, once more
  \item \emph{Say it:} \emph{rúguǒ tiānqì hǎo fēijī jiù yǒuwàng zhèngdiǎn qǐfēi} — If the weather is good, the plane should take off on time.
  \item \emph{Recall:} say the Mandarin for a misfortune, then the Mandarin for to lose one's grip, then say \emph{yǒuwàng} again
\end{itemize}

\begin{tcolorbox}[breakable,colback=teal!4,colframe=teal!35!black,title={Before you move on}]
What does \zh{有望} mean? (\textquotedblleft{}With good hopes of it; likely to happen\textquotedblright{}.) Say it once more.
\end{tcolorbox}
\section[\texorpdfstring{\zh{望} (wàng)}{望 (wàng)}]{\zh{望} — a small piece and \zh{月}, then \zh{王} under them}

\label{lesson:ZH-W532-wang-hope}



\begin{quote}
The word first, then the character inside it.

\begin{itemize}
\raggedright
  \item \emph{Recall:} say \emph{yǒuwàng} — \textbf{R1}, one lesson back, before it decays
\end{itemize}

\begin{quote}
\zh{月} \zh{王}
\end{quote}

\emph{Read it:} each one aloud, with nothing said first — \emph{yuè}, \emph{wáng}
\end{quote}

\subsection*{Script — a small piece and \zh{月}, then \zh{王} under them}
\hlblockfigure{\detokenize{figures/ZH-W532-wang-hope-filmstrip.pdf}}{How \zh{望} is written, stroke by stroke}

\begin{quote}
\zh{望}
\end{quote}

\textbf{Three pieces: two side by side on top, then one across the foot.} Count the pieces, not the strokes.

\textbf{The top left piece, three strokes}: a dot, a bar under it, then a short vertical that turns right along the bottom.

\textbf{The top right piece, \zh{月}, four strokes}, as you wrote it: a stroke down that falls to the left, a bar that turns down and hooks, then two short bars inside.

\textbf{The bottom piece, \zh{王}, four strokes}, as you wrote it: a bar, a second bar, a vertical down through them, then the long bar at the foot.

3 + 4 + 4 = eleven.

\subsection*{Your turn}
\begin{itemize}
\raggedright
  \item \emph{Write it:} \zh{望}, all eleven, with the model in front of you
  \item \emph{Cover:} write it again from memory, and count the pieces, not the strokes
\end{itemize}

\begin{tcolorbox}[breakable,colback=teal!4,colframe=teal!35!black,title={Before you move on}]
How many pieces does \zh{望} have? (\textbf{Three}.) Which two have you written before? (\textbf{\zh{月} and \zh{王}}.)

Source: \href{https://raw.githubusercontent.com/chanind/hanzi-writer-data/68d10a4b21150cae5e1ebbd223eed289cf32d90c/data/\%E6\%9C\%9B.json}{Hanzi Writer Data \zh{望}, pinned stroke paths and medians}.
\end{tcolorbox}
\section[bābude]{bābude — to be dying to; to be eager for (colloquial)}

\label{lesson:ZH-C532-babude}



\begin{quote}
Before the new one: say the Mandarin for to lose one's grip, then the Mandarin for with good hopes of it.
\end{quote}

\subsection*{You'll want to know: bābude}
\emph{bābude} — \textquotedblleft{}to be dying to; to be eager for (colloquial)\textquotedblright{}.

\begin{quote}
\emph{wǒ bābude mǎshàng fàngjià} — I'm dying for the holidays to start right away.
\end{quote}

Its characters come later; for now it is a word for the ear.

Tone: \zh{不} and \zh{得} are both said light here, bābude, so \zh{不} does not change its tone.

\begin{cousinweb}
\emph{bābude} is built from \emph{bā}, \zh{不} 'not' and \zh{得}, all inherited from Middle Chinese, \emph{pae}, \emph{pjut} and \emph{tok} (Guangyun, in Baxter's transcription). The dictionary glosses the first as the name of a region in the west and of a creature; the sense 'to long for' is not given there, and how the three came to mean 'to be dying to' is not settled.
\end{cousinweb}

\begin{grammarlens}[title={What you've built}]
One more word for saying what would happen.
\end{grammarlens}

\subsection*{Your turn}
\begin{itemize}
\raggedright
  \item \emph{Say it:} \emph{bābude}
  \item \emph{Say it:} \emph{bābude}, once more
  \item \emph{Say it:} \emph{wǒ bābude mǎshàng fàngjià} — I'm dying for the holidays to start right away.
  \item \emph{Recall:} say the Mandarin for to lose one's grip, then the Mandarin for with good hopes of it, then say \emph{bābude} again
\end{itemize}

\begin{tcolorbox}[breakable,colback=teal!4,colframe=teal!35!black,title={Before you move on}]
What does bābude mean? (\textquotedblleft{}To be dying to; to be eager for (colloquial)\textquotedblright{}.) Say it once more.
\end{tcolorbox}
\section[\texorpdfstring{\zh{期望} (qīwàng)}{期望 (qīwàng)}]{\zh{期望} (qīwàng) — to expect much of someone; expectations}

\label{lesson:ZH-C532-qiwang}



\begin{quote}
Before the new one: say the Mandarin for with good hopes of it, then the Mandarin for to be dying to.
\end{quote}

\subsection*{You'll want to know: \zh{期望}}
\textbf{\zh{期望}} — \emph{qīwàng} — \textquotedblleft{}to expect much of someone; expectations\textquotedblright{}.

\begin{quote}
\emph{yàoshi nǐ qīwàng tài gāo jiù huì shīwàng} — If you expect too much, you'll be disappointed.
\end{quote}

\begin{cousinweb}
\textbf{\zh{期望}} (\emph{qīwàng}) is built from \zh{期} 'a set time' and \zh{望} 'to look out for', both inherited from Middle Chinese, \emph{gi} and \emph{mjangH} (Guangyun, in Baxter's transcription), which glosses the first 'a time agreed on'.
\end{cousinweb}

\begin{grammarlens}[title={What you've built}]
One more word for saying what would happen.
\end{grammarlens}

\subsection*{Your turn}
\begin{itemize}
\raggedright
  \item \emph{Say it:} \emph{qīwàng}
  \item \emph{Say it:} \emph{qīwàng}, once more
  \item \emph{Say it:} \emph{yàoshi nǐ qīwàng tài gāo jiù huì shīwàng} — If you expect too much, you'll be disappointed.
  \item \emph{Recall:} say the Mandarin for with good hopes of it, then the Mandarin for to be dying to, then say \emph{qīwàng} again
\end{itemize}

\begin{tcolorbox}[breakable,colback=teal!4,colframe=teal!35!black,title={Before you move on}]
What does \zh{期望} mean? (\textquotedblleft{}To expect much of someone; expectations\textquotedblright{}.) Say it once more.
\end{tcolorbox}
\section[\texorpdfstring{\zh{如意} (rúyì)}{如意 (rúyì)}]{\zh{如意} (rúyì) — as one would wish}

\label{lesson:ZH-C532-ruyi}



\begin{quote}
Before the new one: say the Mandarin for to be dying to, then the Mandarin for to expect much of someone.
\end{quote}

\subsection*{You'll want to know: \zh{如意}}
\textbf{\zh{如意}} — \emph{rúyì} — \textquotedblleft{}as one would wish\textquotedblright{}.

\begin{quote}
\emph{rúguǒ gōngzuò bú tài rúyì jiù huàn yí ge} — If your job isn't quite what you hoped, change it.
\end{quote}

\begin{cousinweb}
\textbf{\zh{如意}} (\emph{rúyì}) is built from \zh{如} 'like' and \zh{意} 'what one means', both inherited from Middle Chinese, \emph{nyo} and \emph{'iH} (Guangyun, in Baxter's transcription), which glosses the second 'intent': as one intends.
\end{cousinweb}

\begin{grammarlens}[title={What you've built}]
One more word for saying what would happen.
\end{grammarlens}

\subsection*{Your turn}
\begin{itemize}
\raggedright
  \item \emph{Say it:} \emph{rúyì}
  \item \emph{Say it:} \emph{rúyì}, once more
  \item \emph{Say it:} \emph{rúguǒ gōngzuò bú tài rúyì jiù huàn yí ge} — If your job isn't quite what you hoped, change it.
  \item \emph{Recall:} say the Mandarin for to be dying to, then the Mandarin for to expect much of someone, then say \emph{rúyì} again
\end{itemize}

\begin{tcolorbox}[breakable,colback=teal!4,colframe=teal!35!black,title={Before you move on}]
What does \zh{如意} mean? (\textquotedblleft{}As one would wish\textquotedblright{}.) Say it once more.
\end{tcolorbox}
\section[zhǐwàng]{zhǐwàng — to count on someone}

\label{lesson:ZH-C532-zhiwang}



\begin{quote}
Before the new one: say the Mandarin for to expect much of someone, then the Mandarin for as one would wish it.
\end{quote}

\subsection*{You'll want to know: zhǐwàng}
\emph{zhǐwàng} — \textquotedblleft{}to count on someone\textquotedblright{}.

\begin{quote}
\emph{nǐ búyào zhǐwàng tā bùrán nǐ huì shīwàng} — Don't count on him, or you'll be disappointed.
\end{quote}

Its characters come later; for now it is a word for the ear.

\begin{cousinweb}
\emph{zhǐwàng} is built from \emph{zhǐ} 'to point to' and \zh{望} 'to look out for', both inherited from Middle Chinese, \emph{tsyijX} and \emph{mjangH} (Guangyun, in Baxter's transcription), which glosses the first 'a finger' and 'to show'.
\end{cousinweb}

\begin{grammarlens}[title={What you've built}]
One more word for saying what would happen.
\end{grammarlens}

\subsection*{Your turn}
\begin{itemize}
\raggedright
  \item \emph{Say it:} \emph{zhǐwàng}
  \item \emph{Say it:} \emph{zhǐwàng}, once more
  \item \emph{Say it:} \emph{nǐ búyào zhǐwàng tā bùrán nǐ huì shīwàng} — Don't count on him, or you'll be disappointed.
  \item \emph{Recall:} say the Mandarin for to expect much of someone, then the Mandarin for as one would wish it, then say \emph{zhǐwàng} again
\end{itemize}

\begin{tcolorbox}[breakable,colback=teal!4,colframe=teal!35!black,title={Before you move on}]
What does zhǐwàng mean? (\textquotedblleft{}To count on someone\textquotedblright{}.) Say it once more.
\end{tcolorbox}
